\begin{afterword}
\summaryhead

The post begins with a story. The author asked Marcello, a young olympiad competitor he was considering as an apprentice, how an AI might learn to solve a Rubik's Cube by itself. Marcello said the AI would need ``complexity'' to do certain things. The author told him not to say ``complexity,'' and eventually explained that the word ``doesn't concentrate your probability mass.'' Marcello saw that the word, like ``emergence,'' named a gap without filling it.

The author then generalizes. Complexity is a legitimate concept with mathematical definitions, but using it to explain something you do not understand is a mistake. The mistake is a quick ``skip-over'' that happens below the level of words, and it is common in AI discussions, in academia and in AI startups. The remedy is to notice the feeling of a blank area on your map. The author and Marcello adopted the word ``magic'' as a placeholder that admits an unsolved problem, where ``complexity'' and ``emergence'' hide one.

\responsehead

The idea in this post is old and simple: a word that names a gap does not fill it. Sidney Harris drew it as a cartoon, collected in 1977, in which a blackboard proof reads ``then a miracle occurs,'' and the post uses his caption without credit. What the post adds is a story about its author.

The story does not bear out its lesson. The author's first answer attacks a claim Marcello did not make, that complexity is ``a goal in itself.'' Marcello's objection, that some amount of complexity must be involved, is correct and is never answered. The author's reason for his own position is that the word ``felt like the wrong move,'' and he explains that such feelings cannot be put into words. The dispute ends with a question about whether Marcello has read the author's essay. Then Marcello states the lesson himself, and the post credits the teaching.

The argument, when it comes, retracts the title. Complexity ``is not a useless concept.'' The error happens ``regardless of what causal node goes behind it.'' So banning the word does nothing; the post admits this and then recommends a different word. It says the mistake lives ``below the level of words'' and offers a verbal convention as the fix.

The detector it proposes is circular as described. The mistake is a thought that ``feels like an explanation but isn't.'' The cure is to ``feel which parts of your map are still blank.'' The faculty that produced the error is asked to catch it. No outside check is offered.

The post's evidence is a set of unnamed crowds: people who said the Internet would ``wake up,'' AGI discussants ``left and right,'' academics under ``huge pressure,'' startups whose life plans would collapse. Not one is quoted. The pressure to sweep gaps under the rug is named only in others, though the post shows its author working on the same AI problems.

In short: an old cartoon's point, retold as an anecdote that never answers the student's one real objection, and supported only by unnamed crowds.
\end{afterword}
